\section{One finite coefficient field for a rank-four residual chain}
Fix a prime $\ell$, a coefficient ring $O_c$ in the prescribed algebraic
closure of $\mathbb Q_\ell$, and the valuation ring $V$ in the
algebraically closed field of the given ambient datum. Let $k$ be its residue field. For each admissible $O_d$ write
$k_d$ for its residue field and $j_d:k_d\hookrightarrow k$ for the map
induced by the given local map $O_d\to V$. These maps are compatible
under inclusions of coefficient rings. The ambient hypothesis says that
each element of $V$ agrees modulo its maximal ideal with the image of
some element of some $O_d$.

Let $\Gamma$ be the supplied compact space with group structure and
$\Phi$ its given frame over $O_c$. Write $I$ for the finite set of local
labels. Over a ring $A$, a point $P=(r,n)$ consists of a representation
$r:\Gamma\to M_4(A)$ and one matrix $n_v$ for each $v\in I$. A link
from $P_0$ to $P_1$ is a point over $A[X]$ satisfying the original
residual conditions, with evaluations $P_0$ at zero and $P_1$ at one.
A rank-four chain is a good initial point followed by a finite sequence
of links with good terminal points. Goodness here means the disjunction
$\mathrm{Good4}$ specified below. The polynomial links themselves must
descend.

% BEGIN RCD1
\paragraph{RCD1. A uniform ring for all entries of a finite list of points.}
For a finite list of points over $k$ or $k[X]$ whose representation
entries are locally constant, there is one admissible $O_d$ such that,
for every $O_{c'}\supseteq O_d$, every point in the list has a unique
preimage over $k_{c'}$ or $k_{c'}[X]$, respectively. We may take
$O_d\supseteq O_c$.
\begin{proof}
For $z\in k$, choose a lift to $V$ and use the ambient hypothesis to
find $y\in O_b$ with residue image $z$. For every $O_{c'}\supseteq O_b$,
the residue of the same $y$ in $k_{c'}$ maps to $z$.
Any two coefficient rings lie in a third. Their fraction fields are
finite extensions of $\mathbb Q_\ell$ inside its fixed algebraic
closure. The ring of integers of their finite compositum is an
admissible coefficient ring and contains both by integrality.
Induction gives a common ring for any finite list of residues and
the initial $O_c$.

For each entry of $r$, the open level sets of its locally constant
value cover $\Gamma$. A finite subcover gives finitely many values.
There are sixteen entries and only finitely many label matrices.
In the polynomial case this first gives finitely many \emph{polynomials};
the union of their coefficient sets is finite. Apply the preceding
argument once to all the coefficients, across all the points.

For any larger $O_{c'}$, both $j_{c'}$ and its coefficientwise polynomial
extension are injective. Take the unique preimage of each entry.
Matrix multiplication commutes with these inclusions. Injectivity
reflects $r(1)=1$ and $r(gh)=r(g)r(h)$, giving a representation over
the smaller ring. Together with the preimages of the $n_v$ it gives the
unique point. Its representation entries remain locally constant
because their level sets are the same.
\end{proof}
% END RCD1

% BEGIN RCD2
\paragraph{RCD2. Residual conditions and good vertices descend.}
Let $f:A\hookrightarrow B$ be an injective unital map of commutative
rings, let $a:O_c\to A$ be the coefficient map and let $P$ be an
$A$-point on $\Phi$. If its image satisfies $\Phi.\mathrm{IsRes}$
over $f a$, then $P$ satisfies it over $a$. The same implication holds
for $\mathrm{Good4}$.
\begin{proof}
The symplectic adjoint is given by transposition and fixed signed
changes of indices, so it commutes with $f$. Matrix multiplication,
the similitude equation and every entry of
$\operatorname{adj}(n_v)r(g)n_v$ commute with $f$. In particular so
do the diagonal functions $\psi_{v,i}(g)$.

The pair must have the prescribed multiplier, kill $\Phi.N$, be diagonal
at the distinguished element, and have lower unitriangular symplectic
$n_v$. On $\Phi.G_v$ the conjugated matrices must be upper triangular,
and on $\Phi.\mathrm{Mu}_v$ all their diagonal entries must be one.
The residual data prescribe the diagonals of $r(g)$ and of the local
conjugates by the characters $a(\chi_i)$. They also require one of
two support patterns: global support in $n_P$ together with $n_v$
supported in $q_Q$ at every label, or global support in $n_M$.
Each support set specifies the permitted off-diagonal weights.
Every equality and required zero follows by injectivity, keeping the
given support branch. The continuity clause means continuity into the
discrete quotients by powers of the zero ideal, entry by entry.
Exponent one gives local constancy, which an injection reflects.
Conversely every quotient of a locally constant entry is locally
constant. This proves the full residual predicate.

For $\mathrm{Good4}$, a specified root occurs when some $g,i,j$ of that
weight have $r(g)_{ij}\ne0$. Keep those witnesses. The two extra branches
require global support in $a_A$ and $(n_v)_{10}\ne0$ at each label,
together with their specified occurring roots. The ordinary good
branches use occurrences of roots alone or the condition $\mathrm{InV}$,
which requires global support in $v_V$, local support in $q_Q$, and
$(n_v)_{20}\ne0$. Its last branch also chooses, at each label, an
$s\in\Phi.G_v$ where a specified root character differs from one.
Keep the branch and these same $s$. A zero or equality to one
downstairs would give that equality upstairs. This exhausts the tests.
\end{proof}
% END RCD2

% BEGIN RCD3
\paragraph{RCD3. Descent of the entire directed chain.}
Suppose $s,t$ are points over $k$ with
$\Phi.\mathrm{IsRes}(\mathrm{red}_c,s)$ and
$\mathrm{Chain4}(\Phi,\mathrm{red}_c,s,t)$. There is a coefficient
datum $d$ such that for every $c'$ with $O_c,O_d\subseteq O_{c'}$
there are $s',t'$ over $k_{c'}$ mapping to $s,t$. Over the
coefficient map $a:O_c\hookrightarrow O_{c'}\twoheadrightarrow k_{c'}$,
$s'$ satisfies $\Phi.\mathrm{IsRes}(a,s')$ and
$\mathrm{Chain4}(\Phi,a,s',t')$ holds. This is the original
$\mathrm{ChainDesc4}$.
\begin{proof}
Choose the finite link witnesses $P_1,\ldots,P_m$ for the given
reflexive transitive closure. The residual hypothesis on $s$ and
the residual hypothesis on each \emph{polynomial} $P_j$ give local
constancy of their representation entries by the zero-ideal test in
RCD2. Apply RCD1 to this finite list including $s$. For each allowed
$c'$ let $s'$ and $Q_j$ be their unique preimages. On constants
$j_{c'}a=\mathrm{red}_c$ by compatibility of the maps into $V$.
The corresponding square also commutes after including constants
into the polynomial rings. RCD2 gives the residual predicate for
$s'$ and for every $Q_j$ over its polynomial coefficient map.

Evaluation at zero and one commute with coefficientwise inclusion.
Thus $Q_1(0)=s'$ and $Q_j(1)=Q_{j+1}(0)$ whenever the indicated
links exist. In each case the two points agree after an injection
on both $r$ and every $n_v$. Take $t'=Q_m(1)$ for $m>0$.
RCD2 gives goodness of $s'$ and of each $Q_j(1)$ from the original
witnesses. This is a chain with exactly the original orientation
and length. For $m=0$ take $t'=s'$ and use initial goodness.
The ring $O_d$ was fixed before $c'$, giving the claimed quantifiers.

The Case-III application can therefore enlarge coefficients again
to include a seed and other integral data while retaining the chain
at the final ring. Propagating a condition on components of the
deformation space along it uses the separate rank-four theorem.
\end{proof}
% END RCD3
